\documentclass{ximera}

\begin{document}
	\author{Stitz-Zeager}
	\xmtitle{Answers for Intro Limits}{}

\mfpicnumber{1} \opengraphsfile{AnswersforIntroLimits} % mfpic settings added 
\begin{enumerate}

\item  \begin{multicols}{4} \begin{itemize} \item $\ds{\lim_{x \rightarrow -3^{-}} f(x) = 3}$

\item $\ds{\lim_{x \rightarrow -3^{+}} f(x) = 0}$

\item $\ds{\lim_{x \rightarrow -3} f(x)}$ d.n.e.

\item $f(-3) = 0$

\end{itemize}

\end{multicols}

\bigskip

\begin{multicols}{4}

\begin{itemize}

\item $\ds{\lim_{x \rightarrow 0} f(x) = 9}$

\item  $f(0)$

\item $\ds{\lim_{x \rightarrow 3^{-}} f(x) = 0}$

\item $\ds{\lim_{x \rightarrow 3^{+}} f(x) = -\infty}$
\end{itemize}

\end{multicols}

\bigskip


\item  \begin{enumerate}  \item  If $\ds{\lim_{x \rightarrow a} f(x)}$ exists, say $\ds{\lim_{x \rightarrow a} f(x) = L}$ then the $f(x)$ values approach  $L$  as $x \rightarrow a$ from both directions.  Hence both one-sided limits exist.  In particular,  $\ds{\lim_{x \rightarrow a^{-}} f(x) = L}$ and $\ds{\lim_{x \rightarrow a^{+}} f(x) = L}$.

\item  In Example \ref{limitfromgraphex},  $\ds{\lim_{x \rightarrow -1^{-}} f(x)  = 0}$ and $\ds{\lim_{x \rightarrow -1^{+}} f(x) = 4}$ both exist but $\ds{\lim_{x \rightarrow -1} f(x)}$ does not because the two one-sided limits are not equal.

\end{enumerate}

\item  \begin{multicols}{4} \begin{itemize} \item $\ds{\lim_{x \rightarrow -\infty} g(x) = 1}$

\item $\ds{\lim_{x \rightarrow -2^{-}} g(x) = -\infty}$

\item $\ds{\lim_{x \rightarrow -2^{+}} g(x) = \infty}$

\item  $\ds{\lim_{x \rightarrow \infty} g(x) = -\infty}$

\end{itemize}

\end{multicols}

\bigskip

\begin{multicols}{4}

\begin{itemize}

\item $\ds{\lim_{x \rightarrow 0} g(x) = 2}$

\item $\ds{\lim_{x \rightarrow 2^{-}} g(x) = 1.5}$

\item  $g(2) = 1.5$

\item $\ds{\lim_{x \rightarrow 2^{+}} g(x) = 0}$
\end{itemize}

\end{multicols}

\bigskip

\setcounter{HW}{\value{enumi}}
\end{enumerate}
 


\begin{enumerate}
\setcounter{enumi}{\value{HW}}

\item  $\ds{\lim_{x \rightarrow 2} \frac{2x^2+x-3}{x^2-1} = \frac{7}{3}}$

\medskip
  
\item  $\ds{\lim_{x \rightarrow 1} \frac{2x^2+x-3}{x^2-1} = \frac{5}{2}}$

\medskip

\item   $\ds{\lim_{x \rightarrow -1} \frac{2x^2+x-3}{x^2-1}}$ does not exist

\medskip

\setcounter{HW}{\value{enumi}}
\end{enumerate}



\begin{enumerate}
\setcounter{enumi}{\value{HW}}

\item   $\ds{\lim_{x \rightarrow -1^{+}} \frac{2x^2+x-3}{x^2-1} = \infty}$

\medskip

\item  $\ds{\lim_{x \rightarrow 3} \frac{x^2-3x}{x^2-x-6} = \frac{3}{5}}$


\medskip

\item  $\ds{\lim_{x \rightarrow 3} \frac{x^2-6x}{x^2-6x+9} = -\infty}$

\medskip

\setcounter{HW}{\value{enumi}}
\end{enumerate}




\begin{enumerate}
\setcounter{enumi}{\value{HW}}

\item  $\ds{ \lim_{x \rightarrow 3^{-}} \frac{|3x - x^2| }{x-3} = -3}$ 

\medskip    

\item $\ds{\lim_{x \rightarrow 2} \frac{|6-3x|}{x^2-4x+4} = \infty}$

\medskip

\setcounter{HW}{\value{enumi}}
\end{enumerate}

\begin{enumerate}
\setcounter{enumi}{\value{HW}}

\item  $\ds{\lim_{x \rightarrow 1} \frac{\frac{x}{x-2} +1}{x-1} = -2}$

\medskip

\item $\ds{ \lim_{x \rightarrow 2} \frac{\frac{2x}{x+2}  -1}{x-2}  = \frac{1}{4}}$ 

\medskip
     
\item $\ds{ \lim_{h \rightarrow 0} \frac{\frac{1}{2(x+h) - 1}  - \frac{1}{2x-1} }{h} = -\frac{2}{(2x-1)^2} }$ 

\medskip

\setcounter{HW}{\value{enumi}}
\end{enumerate}

\begin{enumerate}
\setcounter{enumi}{\value{HW}}

 \item  $\ds{ \lim_{x \rightarrow -1} \frac{\sqrt{x+5} - 2 }{x+1}  = \frac{1}{4}}$  
 
 \medskip
 
 \item $\ds{\lim_{h \rightarrow 0} \frac{\sqrt{4+h} - 2}{h}} = \frac{1}{4}$
 
 \medskip
   
 \item  $\ds{ \lim_{h \rightarrow 0} \frac{\sqrt{2x+2h-1} - \sqrt{2x-1} }{h}  = \frac{1}{\sqrt{2x-1}}}$  

\medskip

\setcounter{HW}{\value{enumi}}
\end{enumerate}

\begin{enumerate}
\setcounter{enumi}{\value{HW}}

\item  $\ds{ \lim_{x \rightarrow \infty} \frac{3x-4}{2x+1} = \frac{3}{2}}$

\medskip

\item $\ds{ \lim_{x \rightarrow -\infty}  \frac{1-2x}{x-5} = -2}$ 

\medskip

\item $\ds{ \lim_{x \rightarrow -\infty}\frac{2x-1}{x^2+4} = 0}$

\medskip

\setcounter{HW}{\value{enumi}}
\end{enumerate}

\begin{enumerate}
\setcounter{enumi}{\value{HW}}
\item $\ds{ \lim_{x \rightarrow \infty}\frac{\sqrt{4x^2+x-1}}{1-x} = -2}$

\medskip

\item $\ds{ \lim_{x \rightarrow -\infty}\frac{\sqrt{4x^2+x-1}}{1-x}  = 2}$

\medskip

\item $\ds{ \lim_{x \rightarrow \infty} \frac{x^2+2x+3}{4-x} = -\infty}$

\medskip

\setcounter{HW}{\value{enumi}}
\end{enumerate}


\begin{enumerate}
\setcounter{enumi}{\value{HW}}

\item   \begin{enumerate}

\item  If $-1 < x < 0$, then $\lfloor x \rfloor = -1$. So we can rewrite $f(x) = \frac{x}{\lfloor x \rfloor} = -x$.

\smallskip

\item  Using part (a), we can find  $\ds{ \lim_{x \rightarrow 0^{-}} f(x) =  \lim_{x \rightarrow 0^{-}} -x = 0}$.

\smallskip

\item  If $0 < x < 1$, then $\lfloor x \rfloor =0$, hence $f(x) = \frac{x}{\lfloor x \rfloor}$ does not exist as $x \rightarrow 0^{+}$.

\smallskip

\item  Putting parts (b) and (c) together, we have that   $\ds{ \lim_{x \rightarrow 0} f(x)}$   does not exist.

\smallskip

\end{enumerate}

\item Answers vary.

\item  Answers vary.

\item \begin{enumerate} \item    The function $r(x) = \sqrt{5-x}$ is not continuous at $x = 5$ since $r$ is undefined if $x>5$, so $\ds{\lim_{x \rightarrow 5^{+}} r(x)}$ does not exist. However, $r$ is continuous from the left at $x = 5$ since 

$\ds{\lim_{x \rightarrow 5^{-}} r(x) =  \lim_{x \rightarrow 5^{-}} \sqrt{5-x} = 0 = \sqrt{5-5} = r(5)}$.

\smallskip

\item  The function  $F(x) = \lfloor x \rfloor$ is not continuous at $x = 117$ since $\ds{\lim_{x \rightarrow 117^{-}} \lfloor x \rfloor = 116}$ but $\ds{\lim_{x \rightarrow 117^{+}} \lfloor x \rfloor = 117}$.  However, $F$ is continuous from the right at $x = 117$ since $\ds{\lim_{x \rightarrow 117^{+}} \lfloor x \rfloor = 117 = \lfloor 117 \rfloor = F(117)}$. 

\smallskip

\item If $f$ is continuous at $x=a$, then $\ds{\lim_{x \rightarrow a} f(x) = f(a)}$.  This means  $\ds{\lim_{x \rightarrow a^{-}} f(x) = f(a)}$ and  $\ds{\lim_{x \rightarrow a^{+}} f(x) = f(a)}$, so $f$ is continuous from both directions at $x=a$. The converse is also true since  ff $\ds{\lim_{x \rightarrow a^{-}} f(x) = f(a)}$ and  $\ds{\lim_{x \rightarrow a^{+}} f(x) = f(a)}$, then $\ds{\lim_{x \rightarrow a} f(x) = f(a)}$.

\smallskip

\item  The difference between the scenario here and that in Exercise \ref{twosidedonesidedlimitexistexercise} is that here, we know what each of the one-sided limits are: $f(a)$. They can't be different numbers like they could be in  Example \ref{limitfromgraphex}

\smallskip

\end{enumerate}


\item  We are told nothing of the function on the interval $(-0.000001, 0.000001 )$ so the function has plenty of opportunities to approach $117$.  


\item \begin{enumerate}

\item \begin{itemize} \item  $x^3 > 1000$ for $x > 10$.

\item  $x^3 > 1000000$ for $x > 100$.

\item  $x^3 > 10^{99}$ for $x > 10^{33}$.

\item  $x^3 > N$ for $x > \sqrt[3]{N}$.

\end{itemize}


\item If $x > \sqrt[3]{N}$, then $x^3 > \left(\sqrt[3]{N} \right)^3 = N$.  Per  Definition \ref{infinitelimitsatinfinity}, given $N>0$, choose $M = \sqrt[3]{N}$.  If $x > M$, then $x^3 > M^3 = N$. Hence,  $\ds{\lim_{x \rightarrow \infty} x^3 = \infty}$.

\smallskip

\item Given $N<0$, choose $M = \sqrt[3]{N}$.  If $x< M$, then $x^3 < M^3 = \left(\sqrt[3]{N}\right)^3 = N$.  Hence,   $\ds{\lim_{x \rightarrow -\infty} x^3 = -\infty}$.

\end{enumerate}
 
\end{enumerate}


\end{document}
